% answers_short.tex —— 2008 年·普通高等学校招生考试·天津卷·文科 · 简短答案速查
% 由 tools/make_short_answers.py 自 answers.tex 的 \gksolhead/\gkanswer 提取，
% 并与 2008/2008_数学真题_天津文_答案要点.md 逐题比对后生成，勿手改（改源 MD 后重跑）。

\begin{gknotice}
\textbf{使用说明}：本节只给最终答案便于快速核对，完整推导见解析版。答案为编写者独立推导，\textbf{非官方标准答案}。
\end{gknotice}

\gkpart{选择题}
\begin{center}\begin{tabular}{ccccccccccc}
\toprule
题号 & 1 & 2 & 3 & 4 & 5 & 6 & 7 & 8 & 9 & 10 \\
\midrule
答案 & A & D & A & B & C & C & B & A & D & B \\
\bottomrule
\end{tabular}\end{center}

\gkpart{填空题}
\begin{itemize}[leftmargin=2.6em]
  \item[\textbf{11.}] $10$。
  \item[\textbf{12.}] $10$。
  \item[\textbf{13.}] $12\pi$。
  \item[\textbf{14.}] $8\sqrt2$。
  \item[\textbf{15.}] $x^2+(y+1)^2=18$。
  \item[\textbf{16.}] $432$。
\end{itemize}

\gkpart{解答题}
\begin{itemize}[leftmargin=2.6em]
  \item[\textbf{17.}] （1）$\omega=2$；（2）最大值 $2+\sqrt2$，$x\in\left\{\dfrac{\pi}{16}+\dfrac{k\pi}2\ \middle|\ k\in\mathbf Z\right\}$。
  \item[\textbf{18.}] （1）$p=\dfrac34$；（2）$\dfrac34$；（3）$\dfrac{11}{32}$。
  \item[\textbf{19.}] （1）证明见解析；（2）$\arccos\dfrac2{\sqrt{11}}$；（3）$\arccos\dfrac4{\sqrt{55}}$。
  \item[\textbf{20.}] （1）证明见解析；（2）$a_n=1+\dfrac{q^{n-1}-1}{q-1}$（$q\neq1$），$a_n=n$（$q=1$）；（3）$q=-\sqrt[3]{2}$，证明见解析。
  \item[\textbf{21.}] （1）减区间 $(-\infty,0)$、$\left(\dfrac12,2\right)$，增区间 $\left(0,\dfrac12\right)$、$(2,+\infty)$；（2）$|a|\leqslant\dfrac83$；（3）$b\leqslant-4$。
  \item[\textbf{22.}] （1）$\dfrac{x^2}4-\dfrac{y^2}5=1$；（2）$\left(-\infty,-\dfrac54\right)\cup\left(-\dfrac{\sqrt5}2,0\right)\cup\left(0,\dfrac{\sqrt5}2\right)\cup\left(\dfrac54,+\infty\right)$。
\end{itemize}
